% Lista de siglas e abreviaturas editoriais (curso Set Theory PT)

\clearpage
\phantomsection
\chapter*{Siglas e abreviaturas}
\markboth{\MakeUppercase{Siglas e abreviaturas}}{\MakeUppercase{Siglas e abreviaturas}}
\addcontentsline{toc}{chapter}{Siglas e abreviaturas}

\noindent
Abreviaturas editoriais e siglas de teorias listadas abaixo. Símbolos formais e convenções de escrita estão no capítulo \emph{Notação}; vocabulário matemático em prosa, no \emph{Glossário}.

\begin{description}
\item[Cf.] \emph{Conferir} (latim \textit{confer}): indica ao leitor que deve comparar ou consultar outra passagem.
\item[O.s.] \emph{Ou seja} (equivalente habitual do inglês \textit{i.e.} / \textit{id est}).
\item[p.ex.; P.ex.] \emph{Por exemplo} (equivalente do inglês \textit{e.g.}); \emph{P.ex.} reserva-se ao início de frase.
\item[ca.] \textit{Circa}: por volta de (datas, sobretudo em créditos de imagens).
\item[sse] \emph{Se e somente se} (equivalente do inglês \textit{iff} / \textit{if and only if}).
\item[$\Z$] Teoria de Zermelo \emph{com} o axioma de Fundação.
\item[$\Zminus$] Teoria de Zermelo \emph{sem} o axioma de Fundação.
\item[$\ZF$, $\ZFminus$] Teoria de Zermelo--Fraenkel; a variante $\ZFminus$ dispensa Fundação.
\item[$\ZFC$] Zermelo--Fraenkel com o Axioma da Escolha.
\item[$\LT$] Teoria dos Níveis (\textit{Level Theory}).
\item[$\Zr$] Extensão de $\LT$ com axiomas de Infinito e Conjunto potência, na qual todo o conjunto tem posto.
\item[$\Th{T}$] Letra ou nome curto para uma \emph{teoria} (conjunto de sentenças) de primeira ordem com símbolo primitivo~$\in$.
\item[HC] Hipótese do Contínuo.
\item[HGC] Hipótese Generalizada do Contínuo.
\item[(H)HC] Referência conjunta à HC ou à HGC quando o raciocínio é paralelo.
\item[TS] \emph{Tarski--Scott}: linha de definição de cardinalidade que evita o teorema do bom ordenamento (capítulo correspondente).
\end{description}
